\section{Read black-hole spectra from quantum periods and horizon monodromy}
\label{sec:bh-exact-wkb}
% BEGIN CURATED SUBJECT INDEX
\index{monodromy}
\index{horizon boundary condition}
% END CURATED SUBJECT INDEX

\subsection{Spectral curve and complex horizons}
\label{subsec:bh-spectral-curve}
% BEGIN CURATED SUBJECT INDEX
\index{boundary condition!incoming and outgoing}
\index{resonance!residue}
\index{Stokes curve}
\index{turning point}
\index{monodromy}
\index{quasinormal mode}
\index{horizon boundary condition}
% END CURATED SUBJECT INDEX

Equation~\eqref{eq:bh-master} has the exact-WKB form with
\begin{equation}
  Q(r_*,\omega)=V(r(r_*))-\omega^2 .
  \label{eq:bh-Q}
\end{equation}
Turning points solve $V(r)=\omega^2$.  In the complex $r$ plane, horizons are
singular points of the transformation to $r_*$.  Logarithms generate branch
cuts and Stokes curves that can spiral around a horizon.  An analytic
continuation from the event horizon to infinity must account for these spirals
and for the residues of $\vsub{S}{odd}$ at regular singular points
\cite{Miyachi:2025ptm}.

The QNM condition is again a vanishing incoming coefficient.  Starting with
an ingoing solution at the event horizon, transport it across the complex
Stokes graph and require the incoming component at the outer boundary to
vanish.  This is an open-path connection problem.  Highly damped monodromy
methods instead use closed contours around horizons and singularities.  The
two approaches share connection data but organize approximations differently
\cite{Motl:2003cd}.

\subsection{Quantum periods and Borel--Pad\'e resummation}
\label{subsec:bh-quantum-periods}
% BEGIN CURATED SUBJECT INDEX
\index{Borel transform}
\index{quasinormal mode}
\index{Reissner--Nordstr\"om black hole}
\index{Kerr black hole}
\index{Pad\'e approximant}
% END CURATED SUBJECT INDEX

Let $\gamma_i$ be cycles of the black-hole spectral curve.  The quantum
periods are Borel-resummed Voros integrals
\begin{equation}
  \Pi_i(\omega,\hbar)
  =\mathcal S_{\varphi}
  \oint_{\gamma_i}\vsub{S}{odd}(r,\omega,\hbar)\rmd r .
  \label{eq:bh-quantum-period}
\end{equation}
In black-hole perturbation theory $\hbar$ is often a bookkeeping parameter and
is set to one after the WKB coefficients are generated.  A quantization
condition has the form
\begin{equation}
  \mathcal Q\left(
  \rme^{\Pi_1},\ldots,\rme^{\Pi_g};\omega
  \right)=0 ,
  \label{eq:bh-exact-quantization}
\end{equation}
where $g$ is the number of independent periods relevant to the connection
problem.  High-order coefficients can be Borel transformed and approximated
by Pad\'e rational functions before Laplace integration.  The procedure has
successfully reproduced black-hole QNMs and has been extended to extremal
Reissner--Nordstr\"om and Kerr quantum periods
\cite{Hatsuda:2019eoj,Hatsuda:2023geo,HatsudaShiga:2026}.

\subsection{Relation to complex scaling}
\label{subsec:bh-wkb-csm-relation}
% BEGIN CURATED SUBJECT INDEX
\index{connection coefficient}
\index{resonance!residue}
\index{exact WKB}
\index{complex scaling}
\index{coupled-channel equation}
\index{quasinormal mode!excitation factor}
\index{continuum response}
% END CURATED SUBJECT INDEX

Complex scaling and exact WKB answer different numerical questions about the
same analytic object.  Complex scaling converts the boundary problem into a
matrix eigenproblem and displays the rotated continuum explicitly.  Exact WKB
reduces the pole condition to a small number of resummed periods and
connection multipliers.  The former is well suited to coupled channels and
continuum response; the latter can achieve high precision when the spectral
curve and its cycles are tractable.

Their agreement can be checked at three levels.  Pole locations should match.
The scaling angle should select the same subdominant sectors as the WKB
contour.  Finally, residues computed from left--right eigenvectors should match
derivatives of the exact connection coefficient.  The third comparison is
largely unexplored and would connect ringdown excitation factors directly to
resurgent data.

\subsection{Kerr as the next major test}
\label{subsec:kerr-outlook}
% BEGIN CURATED SUBJECT INDEX
\index{spectrum!point}
\index{GNS construction}
\index{turning point}
\index{exact WKB}
\index{exact quantization condition}
\index{radial equation}
\index{horizon boundary condition}
\index{Kerr black hole}
\index{superradiance}
\index{Teukolsky equation}
% END CURATED SUBJECT INDEX

Kerr perturbations are governed by the Teukolsky equation
\cite{Teukolsky:1972my}.  Separation introduces the azimuthal number $m$, the
spin-weighted spheroidal eigenvalue, and a radial equation whose coefficients
depend on $\omega$.  Superradiance changes flux signs for
$\omega-m\Omega_H$, where $\Omega_H$ is the horizon angular velocity.  A
direct CSM implementation must therefore handle a nonlinear, coupled angular--
radial eigenvalue problem and choose analytic branches consistently.

Exact WKB can work with the associated quantum curve, but turning points and
periods move with both $\omega$ and the angular eigenvalue.  A promising
hybrid strategy is to solve the angular problem spectrally, construct its
analytic derivative with respect to $\omega$, and insert it into either a
nonlinear complex-scaled solver or a resummed radial quantization condition.
Cross-validation against Leaver's method remains essential.

\subsection{Late-time tails and branch cuts}
\label{subsec:bh-tails}
% BEGIN CURATED SUBJECT INDEX
\index{topology}
\index{branch cut}
\index{resonance!residue}
\index{exact WKB}
\index{complex scaling}
\index{quasinormal mode}
\index{late-time tail}
% END CURATED SUBJECT INDEX

In asymptotically flat spacetime, inverse-power tails of $V(r_*)$ produce a
branch point at $\omega=0$.  The late-time signal is governed by the
discontinuity across this cut rather than by the least damped QNM at arbitrarily
late times.  In exact WKB, the threshold appears through coalescing turning
points and a change of Stokes topology.  In complex scaling, it is the endpoint
of the rotated continuum.  A uniform treatment should combine pole residues,
  threshold scaling, and the continuum spectral density rather than attempting
  to approximate the cut by an ever larger set of overtones.
